% GENERATED FILE. Edit canonical lessons, then run npm run generate:books.
% canonical-source-hash: fnv1a64:3b0914650ef61693
% canonical-lessons: LA-C205-vulnus, LA-C205-cicatrix, LA-C205-gravedo, LA-C205-medicamentum, LA-C205-unguentum, LA-R205-first-pass-words-from-chapters-197-200, LA-R205-second-pass-words-from-chapters-201-205

\chapter{Wound, Scar, Head cold, Drug, Ointment}
\label{ch:la-wound-scar-head-cold-drug-ointment}
\hlchaptermodality{205}

\begin{chapteropening}
\textbf{By the end of this chapter:} \emph{I can say a wound, a scar, a head cold, a drug and an ointment.}

Complete the last lesson of chapter 205: I can say a wound, a scar, a head cold, a drug and an ointment.
\end{chapteropening}

\section[vulnus, vulneris]{vulnus, vulneris — a wound}

\label{lesson:LA-C205-vulnus}



\begin{quote}
Before the new one: say the Latin for to throw together; to guess, then the Latin for to cease.
\end{quote}

\subsection*{You'll want to know: vulnus, vulneris}
\textbf{vulnus, vulneris} — \textquotedblleft{}a wound\textquotedblright{}.

\begin{grammarlens}[title={What you've built}]
One more word for this chapter.
\end{grammarlens}

\subsection*{Your turn}
\begin{itemize}
\raggedright
  \item \emph{Say it:} \emph{vulnus, vulneris}
  \item \emph{Say it:} \emph{vulnus, vulneris}, once more
  \item \emph{Say it:} say \emph{vulnus, vulneris}
  \item \emph{Recall:} say the Latin for to throw together; to guess, then the Latin for to cease, then say \emph{vulnus, vulneris} again
\end{itemize}

\begin{tcolorbox}[breakable,colback=teal!4,colframe=teal!35!black,title={Before you move on}]
What does vulnus, vulneris mean? (\textquotedblleft{}A wound\textquotedblright{}.) Say it once more.
\end{tcolorbox}
\section[cicātrīx, cicātrīcis]{cicātrīx, cicātrīcis — a scar}

\label{lesson:LA-C205-cicatrix}



\begin{quote}
Before the new one: say the Latin for to cease, then the Latin for a wound.
\end{quote}

\subsection*{You'll want to know: cicātrīx, cicātrīcis}
\textbf{cicātrīx, cicātrīcis} — \textquotedblleft{}a scar\textquotedblright{}.

\begin{grammarlens}[title={What you've built}]
One more word for this chapter.
\end{grammarlens}

\subsection*{Your turn}
\begin{itemize}
\raggedright
  \item \emph{Say it:} \emph{cicātrīx, cicātrīcis}
  \item \emph{Say it:} \emph{cicātrīx, cicātrīcis}, once more
  \item \emph{Say it:} say \emph{cicātrīx, cicātrīcis}
  \item \emph{Recall:} say the Latin for to cease, then the Latin for a wound, then say \emph{cicātrīx, cicātrīcis} again
\end{itemize}

\begin{tcolorbox}[breakable,colback=teal!4,colframe=teal!35!black,title={Before you move on}]
What does cicātrīx, cicātrīcis mean? (\textquotedblleft{}A scar\textquotedblright{}.) Say it once more.
\end{tcolorbox}
\section[gravēdō, gravēdinis]{gravēdō, gravēdinis — a head cold}

\label{lesson:LA-C205-gravedo}



\begin{quote}
Before the new one: say the Latin for a wound, then the Latin for a scar.
\end{quote}

\subsection*{You'll want to know: gravēdō, gravēdinis}
\textbf{gravēdō, gravēdinis} — \textquotedblleft{}a head cold\textquotedblright{}.

\begin{grammarlens}[title={What you've built}]
One more word for this chapter.
\end{grammarlens}

\subsection*{Your turn}
\begin{itemize}
\raggedright
  \item \emph{Say it:} \emph{gravēdō, gravēdinis}
  \item \emph{Say it:} \emph{gravēdō, gravēdinis}, once more
  \item \emph{Say it:} say \emph{gravēdō, gravēdinis}
  \item \emph{Recall:} say the Latin for a wound, then the Latin for a scar, then say \emph{gravēdō, gravēdinis} again
\end{itemize}

\begin{tcolorbox}[breakable,colback=teal!4,colframe=teal!35!black,title={Before you move on}]
What does gravēdō, gravēdinis mean? (\textquotedblleft{}A head cold\textquotedblright{}.) Say it once more.
\end{tcolorbox}
\section[medicāmentum, medicāmentī]{medicāmentum, medicāmentī — a drug, a medicine}

\label{lesson:LA-C205-medicamentum}



\begin{quote}
Before the new one: say the Latin for a scar, then the Latin for a head cold.
\end{quote}

\subsection*{You'll want to know: medicāmentum, medicāmentī}
\textbf{medicāmentum, medicāmentī} — \textquotedblleft{}a drug, a medicine\textquotedblright{}.

\begin{grammarlens}[title={What you've built}]
One more word for this chapter.
\end{grammarlens}

\subsection*{Your turn}
\begin{itemize}
\raggedright
  \item \emph{Say it:} \emph{medicāmentum, medicāmentī}
  \item \emph{Say it:} \emph{medicāmentum, medicāmentī}, once more
  \item \emph{Say it:} say \emph{medicāmentum, medicāmentī}
  \item \emph{Recall:} say the Latin for a scar, then the Latin for a head cold, then say \emph{medicāmentum, medicāmentī} again
\end{itemize}

\begin{tcolorbox}[breakable,colback=teal!4,colframe=teal!35!black,title={Before you move on}]
What does medicāmentum, medicāmentī mean? (\textquotedblleft{}A drug, a medicine\textquotedblright{}.) Say it once more.
\end{tcolorbox}
\section[unguentum, unguentī]{unguentum, unguentī — an ointment}

\label{lesson:LA-C205-unguentum}



\begin{quote}
Before the new one: say the Latin for a head cold, then the Latin for a drug.
\end{quote}

\subsection*{You'll want to know: unguentum, unguentī}
\textbf{unguentum, unguentī} — \textquotedblleft{}an ointment\textquotedblright{}.

\begin{grammarlens}[title={What you've built}]
One more word for this chapter.
\end{grammarlens}

\subsection*{Your turn}
\begin{itemize}
\raggedright
  \item \emph{Say it:} \emph{unguentum, unguentī}
  \item \emph{Say it:} \emph{unguentum, unguentī}, once more
  \item \emph{Say it:} say \emph{unguentum, unguentī}
  \item \emph{Recall:} say the Latin for a head cold, then the Latin for a drug, then say \emph{unguentum, unguentī} again
\end{itemize}

\begin{tcolorbox}[breakable,colback=teal!4,colframe=teal!35!black,title={Before you move on}]
What does unguentum, unguentī mean? (\textquotedblleft{}An ointment\textquotedblright{}.) Say it once more.
\end{tcolorbox}
\section[(dialogue)]{Words from chapters 197-200 — a first pass}

\label{lesson:LA-R205-first-pass-words-from-chapters-197-200}



\begin{quote}
Say the words for a drug and an ointment.
\end{quote}

\subsection*{Your turn}
Say these aloud:

\begin{itemize}
\raggedright
  \item to kiss, to greet, to invite, to explain, to interpret
  \item to forbid, to defend, to die, to marry, to rest
  \item to compare; to get ready, to persuade, to blame, to share, to avoid
  \item to lift, to reach, to recognise, to refuse, to fly about
\end{itemize}

\begin{tcolorbox}[breakable,colback=teal!4,colframe=teal!35!black,title={Before you move on}]
To kiss? (\textbf{ōsculor, ōsculārī}.) To greet? (\textbf{salūtō, salūtāre}.) To invite? (\textbf{invītō, invītāre}.) To explain? (\textbf{explicō, explicāre}.) To interpret? (\textbf{interpretor, interpretārī}.) To forbid? (\textbf{vetō, vetāre}.) To defend? (\textbf{dēfendō, dēfendere}.) To die? (\textbf{morior, morī}.) To marry? (\textbf{nūbō, nūbere}.) To rest? (\textbf{quiēscō, quiēscere}.) To compare; to get ready? (\textbf{comparō, comparāre}.) To persuade? (\textbf{persuādeō, persuādēre}.) To blame? (\textbf{reprehendō, reprehendere}.) To share? (\textbf{partior, partīrī}.) To avoid? (\textbf{vītō, vītāre}.) To lift? (\textbf{tollō, tollere}.) To reach? (\textbf{perveniō, pervenīre}.) To recognise? (\textbf{agnōscō, agnōscere}.) To refuse? (\textbf{recūsō, recūsāre}.) To fly about? (\textbf{volitō, volitāre}.)
\end{tcolorbox}
\section[(dialogue)]{Words from chapters 201-205 — a second pass}

\label{lesson:LA-R205-second-pass-words-from-chapters-201-205}



\begin{quote}
Say the words for to talk together and to hand over.
\end{quote}

\subsection*{Your turn}
Say these aloud:

\begin{itemize}
\raggedright
  \item to talk together, to hand over, to straighten, to celebrate, to pursue
  \item to advise, to tremble, to welcome, to do, to go forward
  \item to doubt, to look for, to use up, to annoy, to forget
  \item to state firmly, to hurt, to decorate, to throw together; to guess, to cease
  \item a wound, a scar, a head cold, a drug, an ointment
\end{itemize}

\begin{tcolorbox}[breakable,colback=teal!4,colframe=teal!35!black,title={Before you move on}]
To talk together? (\textbf{colloquor, colloquī}.) To hand over? (\textbf{trādō, trādere}.) To straighten? (\textbf{corrigō, corrigere}.) To celebrate? (\textbf{celebrō, celebrāre}.) To pursue? (\textbf{persequor, persequī}.) To advise? (\textbf{suādeō, suādēre}.) To tremble? (\textbf{tremō, tremere}.) To welcome? (\textbf{excipiō, excipere}.) To do? (\textbf{agō, agere}.) To go forward? (\textbf{prōcēdō, prōcēdere}.) To doubt? (\textbf{dubitō, dubitāre}.) To look for? (\textbf{quaerō, quaerere}.) To use up? (\textbf{cōnsūmō, cōnsūmere}.) To annoy? (\textbf{vexō, vexāre}.) To forget? (\textbf{oblīvīscor, oblīvīscī}.) To state firmly? (\textbf{affirmō, affirmāre}.) To hurt? (\textbf{doleō, dolēre}.) To decorate? (\textbf{ōrnō, ōrnāre}.) To throw together; to guess? (\textbf{coniciō, conicere}.) To cease? (\textbf{dēsinō, dēsinere}.) A wound? (\textbf{vulnus, vulneris}.) A scar? (\textbf{cicātrīx, cicātrīcis}.) A head cold? (\textbf{gravēdō, gravēdinis}.) A drug? (\textbf{medicāmentum, medicāmentī}.) An ointment? (\textbf{unguentum, unguentī}.)
\end{tcolorbox}
